% =========================================================================
% ALGORITMOS DE CAMINHOS AUMENTANTES
% =========================================================================
\chapter{Algoritmos de Caminhos Aumentantes (Abordagem Primal)}\label{sec:caminhos_aumentantes}

\section{Arquitetura Orientada a Objetos e Estruturas de Dados}

A infraestrutura de algoritmos de fluxo desenvolvida neste projeto (Apêndices~\ref{ap:flownetwork} a~\ref{ap:push_relabel_improved}) é implementada em C++23 fundamentada em orientação a objetos.

Definem-se apelidos de tipos inteiros padronizados (\lstinline{Long} e \lstinline{Size}) e constantes de saturação (\lstinline{INF} e \lstinline{MAX} com deslocamento de bits \lstinline{>> 8} para prevenir estouro aritmético em operações de soma).

O núcleo do sistema é definido pela classe base abstrata \lstinline{FlowNetwork}, que encapsula a representação do digrafo e o estado do fluxo:

\begin{itemize}
    \item \textbf{\lstinline{Edge}:} Estrutura que representa um arco direcionado, contendo \lstinline{from} (origem), \lstinline{to} (destino), \lstinline{capacity} (capacidade original $c$) e \lstinline{flow} (fluxo atual $f$).
    \item \textbf{\lstinline{size}:} Quantidade de vértices ($|V|$), utilizada no dimensionamento de estruturas auxiliares.
    \item \textbf{\lstinline{edges}:} Vetor contíguo que armazena todos os arcos da rede.
    \item \textbf{\lstinline{adjacency}:} Vetor de listas de adjacência (\lstinline{std::vector<std::vector<Size>>}), armazenando os índices dos arcos no vetor \lstinline{edges} para acesso em tempo $\mathcal{O}(1)$.
\end{itemize}

A manipulação de fluxo é encapsulada em dois métodos protegidos:

\begin{itemize}
    \item \textbf{\lstinline{get_residual_capacity(Size edge_id)}:} Computa a capacidade residual $c_f(u,v) = c(u,v) - f(u,v)$. O atributo \lstinline{[[nodiscard]]} assegura que o valor retornado seja consumido pelo chamador.
    \item \textbf{\lstinline{push_flow(Size edge_id, Long flow_amount)}:} Atualiza o fluxo no arco direto e no reverso. Como arcos direto e reverso são inseridos em pares consecutivos (índices $2k$ e $2k+1$), a operação bit a bit \lstinline{id ^ 1ULL} acessa o arco complementar em $\mathcal{O}(1)$.
\end{itemize}

A interface pública fornece métodos virtuais \lstinline{make} e \lstinline{clone} para instanciação e clonagem profunda de redes sem acoplamento a classes concretas. O método virtual puro \lstinline{compute_max_flow(Size source, Size sink)} padroniza a execução do cálculo de fluxo máximo, e \lstinline{get_edges()} permite consultar os arcos resultantes via referência constante (\lstinline{const &}). A implementação completa da classe base está no Código~\ref{lst:flownetwork} (Apêndice~\ref{ap:flownetwork}).


\section{O Paradigma dos Caminhos Aumentantes}
\label{sec:fluxo_primal}

Algoritmos primais mantêm a viabilidade das restrições de conservação ($\sum f_{in} = \sum f_{out}$) e capacidade ($0 \le f \le c$) em todas as iterações.

Partindo do fluxo nulo ($f \equiv 0$), o algoritmo busca na rede residual $D_f$ um caminho simples direcionado da fonte $s$ ao sumidouro $t$ (caminho aumentante). Identificado o caminho, aumenta-se o fluxo pelo valor da capacidade residual mínima ao longo dos seus arcos (capacidade gargalo). A condição de otimalidade é atingida quando não existem caminhos aumentantes conectando $s$ a $t$ em $D_f$.

A estratégia de seleção dos caminhos determina a complexidade assintótica dos algoritmos:
\begin{itemize}
    \item \textbf{Ford-Fulkerson:} Seleção arbitrária via busca em profundidade (DFS); tempo pseudopolinomial $\mathcal{O}(|A| \cdot |f^*|)$.
    \item \textbf{Edmonds-Karp:} Seleção de caminhos com menor número de arcos via busca em largura (BFS); tempo fortemente polinomial $\mathcal{O}(|V| \cdot |A|^2)$.
    \item \textbf{Dinic:} Busca em fases sobre digrafos de níveis com fluxos bloqueadores; tempo $\mathcal{O}(|V|^2 \cdot |A|)$.
\end{itemize}


\section{O Método Genérico de Ford-Fulkerson e a Barreira Pseudopolinomial}\label{sec:ford_fulkerson}

\paragraph{Princípio Teórico.}
O método clássico de Ford e Fulkerson~\cite{ford1956maximal} baseia-se na manutenção da rede residual $D_f$ e no aumento sucessivo de fluxo ao longo de caminhos aumentantes. O método original não especifica a ordem de escolha dos caminhos em digrafos com múltiplos trajetos viáveis~\cite{ahuja1993network}.

\paragraph{Complexidade Assintótica e Barreira Pseudopolinomial.}
Caso as capacidades assumam valores irracionais, seleções arbitrárias de caminhos podem convergir para um valor estritamente inferior ao fluxo máximo~\cite{cook1998combinatorial}.

Para capacidades inteiras, uma sequência desfavorável de caminhos pode demandar tempo $\mathcal{O}(|A| \cdot |f^*|)$, dependente da magnitude das capacidades originais. Essa dependência caracteriza a barreira pseudopolinomial, ilustrada no contraexemplo da Figura~\ref{fig:ford_fulkerson_patologico}.

\begin{figure}[!ht]
\centering
\begin{tikzpicture}[thick]
\begin{scope}[xshift=0cm]
  \tikzsubcaption{(2.5,2.6)}{(a) Rede patológica ($M = 10^6$)}
  \node[source node] (s) at (0, 1) {$s$};
  \node[flow node] (u) at (2.5, 2) {$u$};
  \node[flow node] (v) at (2.5, 0) {$v$};
  \node[sink node] (t) at (5, 1) {$t$};

  \draw[flow edge] (s) -- node[edge lbl above] {$M$} (u);
  \draw[flow edge] (s) -- node[edge lbl below] {$M$} (v);
  \draw[sat edge] (u) -- node[edge lbl right, text=red!80!black] {$c=1$} (v);
  \draw[flow edge] (u) -- node[edge lbl above] {$M$} (t);
  \draw[flow edge] (v) -- node[edge lbl below] {$M$} (t);
\end{scope}

\begin{scope}[xshift=7.2cm]
  \tikzsubcaption{(2.5,2.6)}{(b) Degeneração da DFS: $2M$ iterações}
  \node[source node] (s2) at (0, 1) {$s$};
  \node[flow node] (u2) at (2.5, 2) {$u$};
  \node[flow node] (v2) at (2.5, 0) {$v$};
  \node[sink node] (t2) at (5, 1) {$t$};

  \draw[tree edge, line width=1.8pt] (s2) -- node[edge lbl above] {$+1$} (u2);
  \draw[sat edge, line width=1.8pt] (u2) -- node[edge lbl right] {$+1$} (v2);
  \draw[tree edge, line width=1.8pt] (v2) -- node[edge lbl below] {$+1$} (t2);

  \draw[flow edge, opacity=0.3] (s2) -- (v2);
  \draw[flow edge, opacity=0.3] (u2) -- (t2);

  \node[font=\scriptsize, align=center, below=3pt] at (2.5,-0.5) {DFS alterna $s \to u \to v \to t$ e $s \to v \to u \to t$\\Aumento $\delta=1 \implies 2 \times 10^6$ passos (BFS / Dinic: 1 fase / 2 passos)};
\end{scope}
\end{tikzpicture}

\caption{Contraexemplo patológico de Ford-Fulkerson com capacidades $M = 10^6$ e arco central unitário $c=1$. (a)~Rede com fluxo máximo $|f^*| = 2M$. (b)~DFS ingênua alternando caminhos pelo arco central, aumentando $\delta=1$ por passo em $2 \times 10^6$ iterações. Em contrapartida, métodos baseados em BFS priorizam caminhos de menor distância: Edmonds-Karp encerra em $2$ iterações e Dinic em $1$ única fase de fluxo bloqueador.}
\label{fig:ford_fulkerson_patologico}
\end{figure}

\paragraph{Implementação em C++.}
A implementação do método com DFS está descrita no Código~\ref{lst:ford_fulkerson} (Apêndice~\ref{ap:ford_fulkerson}).


\section{O Algoritmo de Edmonds-Karp}\label{sec:edmonds_karp}

\paragraph{Princípio Teórico e Seleção por Caminhos Mais Curtos.}
O algoritmo de Edmonds-Karp~\cite{edmonds1972theoretical, dinic1970algorithm} seleciona a cada iteração o caminho aumentante com menor número de arcos entre $s$ e $t$ na rede residual $D_f$. Cada arco recebe peso unitário para fins de distância pela BFS.

\paragraph{Dinâmica Operacional via BFS.}
A BFS garante a monotonicidade das distâncias: ao longo das iterações, a distância mínima em número de arcos $d(s, v)$ em $D_f$ não decresce para nenhum vértice $v \in V$~\cite{cook1998combinatorial}.

\paragraph{Complexidade Assintótica.}
Quando um arco $(u,v)$ atua como gargalo, ele é saturado e removido de $D_f$, satisfazendo $d(v) = d(u) + 1$. Para que retorne a $D_f$ em iteração posterior, fluxo deve ser enviado em sentido oposto $(v,u)$, exigindo que a nova distância satisfaça $d'(u) = d'(v) + 1$. Pela monotonicidade ($d'(v) \ge d(v)$), decorre:
\begin{equation}
    d'(u) \ge d(u) + 2
\end{equation}
Como o comprimento de um caminho simples é limitado por $|V| - 1$, cada arco pode ser saturado no máximo $\mathcal{O}(|V|)$ vezes. Havendo $\mathcal{O}(|A|)$ arcos, o número total de aumentos é $\mathcal{O}(|V| \cdot |A|)$. Com custo $\mathcal{O}(|A|)$ por BFS, o tempo total no pior caso é $\mathcal{O}(|V| \cdot |A|^2)$~\cite{korte2018combinatorial}.

\paragraph{Implementação em C++.}
A classe correspondente está estruturada como \lstinline{EdmondsKarp} (Código~\ref{lst:edmonds_karp}, Apêndice~\ref{ap:edmonds_karp}).

\begin{figure}[!ht]
\centering
\begin{tikzpicture}[thick]
  \node[source node] (s) at (0, 0) {$s$ (0)};
  \node[flow node] (a) at (3.5, 1.5) {$a$ (1)};
  \node[flow node] (b) at (2, -1.5) {$b$ (1)};
  \node[flow node] (c) at (5, -1.5) {$c$ (2)};
  \node[sink node] (t) at (7, 0) {$t$ (2)};

  \draw[res edge, line width=1.5pt] (s) -- node[edge lbl above] {$c=5$} (a);
  \draw[res edge, line width=1.5pt] (a) -- node[edge lbl above] {$c=5$} (t);

  \draw[flow edge] (s) -- node[edge lbl below] {$c=10$} (b);
  \draw[flow edge] (b) -- node[edge lbl below] {$c=10$} (c);
  \draw[flow edge] (c) -- node[edge lbl below] {$c=10$} (t);

  \draw[flow edge] (a) -- node[edge lbl right] {$c=2$} (b);
\end{tikzpicture}
\caption{O Algoritmo de Edmonds-Karp utiliza busca em largura (BFS) na rede residual para encontrar o caminho aumentante mais curto em número de arestas. A BFS seleciona o caminho $s \to a \to t$ (comprimento 2), preterindo $s \to b \to c \to t$ (comprimento 3), mesmo com capacidade superior. Os níveis estão indicados entre parênteses.}
\label{fig:edmonds_karp}
\end{figure}


\section{O Algoritmo de Dinic}\label{sec:dinic}

\paragraph{Princípio Teórico e Digrafo de Níveis.}
O algoritmo de Dinic~\cite{dinic1970algorithm} reduz o tempo assintótico ao processar o fluxo em fases estruturadas sobre um \textbf{Digrafo de Níveis} (\textit{Level Graph}).

No início de cada fase, uma BFS a partir de $s$ calcula as distâncias mínimas $d(v)$ em $D_f$. O Digrafo de Níveis é o subgrafo acíclico composto exclusivamente pelos arcos residuais que satisfazem:
\begin{equation}
    d(v) = d(u) + 1
\end{equation}

\paragraph{Fluxos Bloqueadores.}
Sobre o Digrafo de Níveis, emprega-se busca em profundidade (DFS) para encontrar um \textbf{Fluxo Bloqueador} (\textit{Blocking Flow}), saturando simultaneamente caminhos admissíveis até que não reste caminho de $s$ a $t$ no Digrafo de Níveis~\cite{cook1998combinatorial}.

\paragraph{Otimização por Ponteiros de Avanço.}
Para evitar revisitações desnecessárias a arcos saturados ou caminhos sem saída em uma mesma fase, a implementação mantém um vetor de ponteiros (\textit{dead-end pointers}) que registra o índice do último arco explorado a partir de cada vértice.

\paragraph{Complexidade Assintótica.}
A determinação do fluxo bloqueador opera em $\mathcal{O}(|V| \cdot |A|)$ com a otimização de ponteiros~\cite{korte2018combinatorial}. Pela monotonicidade estrita das distâncias entre fases ($d_{k+1}(t) \ge d_k(t) + 1$), o número máximo de fases é limitado por $|V| - 1$. Consequentemente, o tempo total é $\mathcal{O}(|V|^2 \cdot |A|)$~\cite{dinic1970algorithm}.

\paragraph{Implementação em C++.}
A classe correspondente está estruturada como \lstinline{Dinic} (Código~\ref{lst:dinic}, Apêndice~\ref{ap:dinic}).

\begin{figure}[!ht]
\centering
\begin{tikzpicture}[thick]
\foreach \x in {0,2.5,5.0,7.5} { \draw[layer line] (\x,-0.7) -- (\x,3.8); }

\node[gray,font=\footnotesize] at (0,4.1) {nív. 0};
\node[gray,font=\footnotesize] at (2.5,4.1) {nív. 1};
\node[gray,font=\footnotesize] at (5.0,4.1) {nív. 2};
\node[gray,font=\footnotesize] at (7.5,4.1) {nív. 3};

\node[source node] (s) at (0,1.5) {$s$};
\node[flow node] (a) at (2.5,3) {$a$};
\node[flow node] (b) at (2.5,0) {$b$};
\node[flow node] (c) at (5.0,3) {$c$};
\node[flow node] (d) at (5.0,0) {$d$};
\node[sink node] (t) at (7.5,1.5) {$t$};

\draw[flow edge,blue,line width=1.5pt] (s) -- (a);
\draw[flow edge,blue,line width=1.5pt] (s) -- (b);
\draw[flow edge,blue,line width=1.5pt] (a) -- (c);
\draw[flow edge,blue,line width=1.5pt] (b) -- (d);
\draw[flow edge,blue,line width=1.5pt] (a) -- (d);
\draw[flow edge,blue,line width=1.5pt] (b) -- (c);
\draw[flow edge,blue,line width=1.5pt] (c) -- (t);
\draw[flow edge,blue,line width=1.5pt] (d) -- (t);

\draw[nontree edge,bend right=30] (a) to (b);
\draw[nontree edge,bend right=20] (c) to (a);
\end{tikzpicture}

\caption{Digrafo de Níveis no Algoritmo de Dinic (organizado em colunas, onde cada coluna da esquerda para a direita representa um nível crescente). Arcos azuis são os arcos admissíveis (avançam exatamente um nível: $d(v) = d(u)+1$), que compõem o Digrafo de Níveis. Arcos cinza tracejados são excluídos por não avançarem nível. O Fluxo Bloqueador é extraído apenas sobre os arcos azuis.}
\label{fig:grafo_niveis}
\end{figure}


